\documentclass[11pt]{article}
\usepackage[margin=1in]{geometry}
\usepackage{amsmath,amssymb,amsthm}
\usepackage{booktabs}
\usepackage{hyperref}

\theoremstyle{plain}
\newtheorem{thm}{Theorem}
\newtheorem{lem}[thm]{Lemma}
\newtheorem{cor}[thm]{Corollary}
\newtheorem{prop}[thm]{Proposition}
\theoremstyle{definition}
\newtheorem{defn}[thm]{Definition}
\theoremstyle{remark}
\newtheorem{rem}[thm]{Remark}
\newtheorem{conj}[thm]{Conjecture}
\theoremstyle{plain}
\newtheorem{gapthm}{Gap}

\newcommand{\bI}{\mathbf{I}}
\newcommand{\bM}{\mathbf{M}}
\newcommand{\bK}{\mathbf{K}}
\newcommand{\bL}{\mathbf{L}}
\newcommand{\bu}{\mathbf{u}}
\newcommand{\bw}{\mathbf{w}}
\newcommand{\bv}{\mathbf{v}}
\newcommand{\WCt}{\widetilde{WC}}
\newcommand{\At}{\widetilde{A}}
\newcommand{\ord}{\operatorname{ord}}

\title{Collision divisors and the Laurent locus\\[2mm]
\large The jump criterion for Motegi's type $C$ structure constants, proved;\\
and the determinant/Pfaffian stratification shown not to be it}
\author{Clio}
\date{16 October 2026}

\makeatletter
%% protocol-header.tex -- PROTOCOL §2.3 page-1 block, as a TEMPLATE not an instance.
%%
%% PROTOCOL §2.3: "Every PDF must state, on its first page: author, date,
%% recipient, a title, and the commit hash of the work-in-progress push it
%% corresponds to.  A PDF with no commit hash cannot be checked against the
%% record and should be treated as a draft opinion, not a result."
%%
%% Written 2026-10-16 because the fault had recurred twice and I had twice
%% fixed only the instance:
%%   - 2026-10-14: the Q473 PDF went to Rick with no recipient and no hash.
%%   - 2026-10-14: the Q460 PDF went to Rick saying "Recipient: Robin Langer".
%%     He noticed, and took the commit hash as authority instead.
%% Both times I repaired the one file and wrote "fix the template, not just the
%% instance" in my notes.  This file is that fix.  \input it; do not retype it.
%%
%% USAGE, immediately after \maketitle:
%%   %% protocol-header.tex -- PROTOCOL §2.3 page-1 block, as a TEMPLATE not an instance.
%%
%% PROTOCOL §2.3: "Every PDF must state, on its first page: author, date,
%% recipient, a title, and the commit hash of the work-in-progress push it
%% corresponds to.  A PDF with no commit hash cannot be checked against the
%% record and should be treated as a draft opinion, not a result."
%%
%% Written 2026-10-16 because the fault had recurred twice and I had twice
%% fixed only the instance:
%%   - 2026-10-14: the Q473 PDF went to Rick with no recipient and no hash.
%%   - 2026-10-14: the Q460 PDF went to Rick saying "Recipient: Robin Langer".
%%     He noticed, and took the commit hash as authority instead.
%% Both times I repaired the one file and wrote "fix the template, not just the
%% instance" in my notes.  This file is that fix.  \input it; do not retype it.
%%
%% USAGE, immediately after \maketitle:
%%   %% protocol-header.tex -- PROTOCOL §2.3 page-1 block, as a TEMPLATE not an instance.
%%
%% PROTOCOL §2.3: "Every PDF must state, on its first page: author, date,
%% recipient, a title, and the commit hash of the work-in-progress push it
%% corresponds to.  A PDF with no commit hash cannot be checked against the
%% record and should be treated as a draft opinion, not a result."
%%
%% Written 2026-10-16 because the fault had recurred twice and I had twice
%% fixed only the instance:
%%   - 2026-10-14: the Q473 PDF went to Rick with no recipient and no hash.
%%   - 2026-10-14: the Q460 PDF went to Rick saying "Recipient: Robin Langer".
%%     He noticed, and took the commit hash as authority instead.
%% Both times I repaired the one file and wrote "fix the template, not just the
%% instance" in my notes.  This file is that fix.  \input it; do not retype it.
%%
%% USAGE, immediately after \maketitle:
%%   \input{protocol-header}
%%   \protocolheader{Rick (grandpa-rick)}{clio-vega/proofs}{60e4ce3}
%%
%% The recipient argument is MANDATORY and has no default on purpose: a default
%% is what produced "Recipient: Robin Langer" on a PDF addressed to Rick.

%% \maketitle CLEARS \@date (and \@author, \@title) in the article class, so a
%% macro called AFTER \maketitle that reads \@date prints an EMPTY row.  Caught
%% 2026-10-16 only because I ran pdftotext on page 1 instead of trusting rc=0:
%% the build was clean and the Date field was blank.  Hence: snapshot the date
%% HERE, at \input time in the preamble, where \@date is still live.
\let\protocol@date\@date

\newcommand{\protocolheader}[3]{%
  \begin{center}
  \fbox{\begin{minipage}{0.92\textwidth}
  \small
  \textbf{PROTOCOL §2.3 record.}\\[1mm]
  \begin{tabular}{@{}ll@{}}
  Author     & Clio \\
  Date       & \protocol@date \\
  Recipient  & #1 \\
  Repository & \texttt{#2} \\
  Commit     & \texttt{#3} \\
  \end{tabular}\\[1mm]
  This PDF corresponds to the pushed commit above.  Where this text and the
  commit disagree, \textbf{the commit is authoritative}.
  \end{minipage}}
  \end{center}
  \vspace{2mm}%
}

%%   \protocolheader{Rick (grandpa-rick)}{clio-vega/proofs}{60e4ce3}
%%
%% The recipient argument is MANDATORY and has no default on purpose: a default
%% is what produced "Recipient: Robin Langer" on a PDF addressed to Rick.

%% \maketitle CLEARS \@date (and \@author, \@title) in the article class, so a
%% macro called AFTER \maketitle that reads \@date prints an EMPTY row.  Caught
%% 2026-10-16 only because I ran pdftotext on page 1 instead of trusting rc=0:
%% the build was clean and the Date field was blank.  Hence: snapshot the date
%% HERE, at \input time in the preamble, where \@date is still live.
\let\protocol@date\@date

\newcommand{\protocolheader}[3]{%
  \begin{center}
  \fbox{\begin{minipage}{0.92\textwidth}
  \small
  \textbf{PROTOCOL §2.3 record.}\\[1mm]
  \begin{tabular}{@{}ll@{}}
  Author     & Clio \\
  Date       & \protocol@date \\
  Recipient  & #1 \\
  Repository & \texttt{#2} \\
  Commit     & \texttt{#3} \\
  \end{tabular}\\[1mm]
  This PDF corresponds to the pushed commit above.  Where this text and the
  commit disagree, \textbf{the commit is authoritative}.
  \end{minipage}}
  \end{center}
  \vspace{2mm}%
}

%%   \protocolheader{Rick (grandpa-rick)}{clio-vega/proofs}{60e4ce3}
%%
%% The recipient argument is MANDATORY and has no default on purpose: a default
%% is what produced "Recipient: Robin Langer" on a PDF addressed to Rick.

%% \maketitle CLEARS \@date (and \@author, \@title) in the article class, so a
%% macro called AFTER \maketitle that reads \@date prints an EMPTY row.  Caught
%% 2026-10-16 only because I ran pdftotext on page 1 instead of trusting rc=0:
%% the build was clean and the Date field was blank.  Hence: snapshot the date
%% HERE, at \input time in the preamble, where \@date is still live.
\let\protocol@date\@date

\newcommand{\protocolheader}[3]{%
  \begin{center}
  \fbox{\begin{minipage}{0.92\textwidth}
  \small
  \textbf{PROTOCOL §2.3 record.}\\[1mm]
  \begin{tabular}{@{}ll@{}}
  Author     & Clio \\
  Date       & \protocol@date \\
  Recipient  & #1 \\
  Repository & \texttt{#2} \\
  Commit     & \texttt{#3} \\
  \end{tabular}\\[1mm]
  This PDF corresponds to the pushed commit above.  Where this text and the
  commit disagree, \textbf{the commit is authoritative}.
  \end{minipage}}
  \end{center}
  \vspace{2mm}%
}

\makeatother

\begin{document}
\maketitle
\protocolheader{Robin (\texttt{langer.robin@gmail.com})}{clio-vega/proofs}{c5f51b8}

\begin{abstract}
Two results, one negative and one positive, about the same object: the
structure constants $c^{\nu}_{\lambda\mu}$ of the fixed-point algebra of
Motegi's type $C$ trigonometric weight functions (arXiv:2610.15344).

\textbf{(1)} Q492 asked whether the determinant/Pfaffian dichotomy for type
$B$/$C$ boundary partition functions is the \emph{same} stratification as the
Laurent criterion $|\nu|-\max(|\lambda|,|\mu|)\le1$ measured in Q484.  It is
not, and the reason is structural rather than numerical.  The integer
governing Motegi's boundary degeneration is $n-s$, not $n-r$; the number $n$ is
a gauge parameter annihilated by his own stability relation; and the
gauge-invariant residue $s-r=|\bM|-|\bI|$ is a cardinality gap between an
\emph{evaluation point} and a \emph{basis label}, whereas the jump is a gap
among \emph{three basis labels}.  The two integers are gaps in different slots
of the same matrix.  Q492 is closed as a negative (Theorem~\ref{thm:gate}).

\textbf{(2)} The criterion itself, which Q484 could only measure ($225$
triples, zero error in either direction), is \emph{proved} in one direction,
and the mechanism is identified.  The candidate pole divisors of
$c^{\nu}_{\lambda\mu}$ split into \emph{collision divisors} --- loci on which
two of the $2^N$ specialisation points literally collide --- and the
\emph{long-root divisors} $w_aw_b=1$, which are not collisions.  On a collision
divisor a row collapse forces the numerator of the triangular recursion to
vanish to matching order, so no pole can occur, for \emph{any}
$(\lambda,\mu,\nu)$ (Theorem~\ref{thm:collision}).  On the long-root divisors a
pole is created exactly when a certain two-element deletion stays above both
$\lambda$ and $\mu$, and that condition is equivalent --- by a two-line
cardinality argument --- to jump $\ge2$ (Lemma~\ref{lem:comb}).  Hence
jump $\le1$ implies Laurent (Theorem~\ref{thm:main}).  The converse requires a
non-cancellation statement which is measured, not proved, and is stated as
Gap~\ref{gap:nonvanishing}.

The mechanism predicts more than the criterion: it predicts the \emph{exact
pole divisor} of each structure constant, hyperplane by hyperplane.  Verified
$20/20$ at $N=2$ and $115/115$ at $N=3$ against SymPy, by an independent
combinatorial mechanism; both pre-registered plants fire.
\end{abstract}

\tableofcontents

\section{The gate: Q492 is a one-coordinate match}
\label{sec:gate}

The session brief armed a thirty-minute gate before any machinery, because two
integers were in play and they shared a noun.  The gate ran, and it closed the
question.

\subsection{What Motegi's integer actually is}

Fix $N$ and let $\bw=(w_1,\dots,w_N)$.  Motegi's partition functions
$F_{N,n,r}(\bu|\bw|\bI)$, and their type $C$ and $B$ specialisations
$WC_{N,n,r}$, $WB_{N,n,r}$, carry three integers (L331--L470 of the e-print
source \texttt{typebcweightfunctions\_arxivversionfinal-3.tex}):
\begin{itemize}
\item $N$, the number of columns, carrying $\bw$;
\item $n$, the number of rows, carrying $\bu=(u_1,\dots,u_n)$;
\item $r=|\bI|$, the number of $1$'s on the bottom boundary, where
  $\bI=(i_1<\dots<i_r)$ equivalently encodes a strict partition
  $N\ge\lambda_1>\dots>\lambda_r\ge1$ by $i_j=\lambda_{r+1-j}$.
\end{itemize}
The specialisation is indexed by a \emph{fourth} datum
$\bM=\{m_1<\dots<m_s\}\subseteq[N]$, with
\begin{equation}
\bu=\bw_{\bM}:=(w_{m_1},\dots,w_{m_s},\underbrace{1,\dots,1}_{n-s}).
\label{eq:spec}
\end{equation}
Motegi states the distinction himself, at L2790--L2803:
\begin{quote}
``Thus $r$ specifies the label cardinality, whereas $s$ specifies the number of
selected $w$-variables in the specialization.''
\end{quote}
So the number of spectral parameters set to $1$ is $n-s$, \textbf{not} $n-r$.

My own wake note had recorded $n-r$, quoting L2094: ``upon setting
$u_{r+1}=\dots=u_n=1$, the configuration in the top $(n-r)$ rows is uniquely
determined by the condition $[K(1)]^1_0=0$.''  That quotation is accurate and
the reading of it was wrong: the proposition it belongs to is the
\emph{equal-cardinality} case $s=r$, where $\bw_{\bM}=(w_{m_1},\dots,w_{m_r},
1,\dots,1)$, so that $n-s$ and $n-r$ coincide.  The quantity was read off a
stratum on which the two integers agree and then carried off it.

\subsection{And $n$ is gauge}

Worse for the hypothesis: $n$ is not an invariant of anything.  Motegi's
stability relation (\texttt{stabilityWC}, L2641--L2647) reads
\begin{equation}
\WCt_{N,n,r}(\bu|\bw|\bI)\big|_{u_n=1}
=\WCt_{N,n-1,r}(\bu\setminus\{u_n\}|\bw|\bI),
\qquad r\le n-1,
\label{eq:stab}
\end{equation}
and iterating it (his \texttt{eq:typeC-stability-s-ge-r}, L2823--L2830) gives,
for $s\ge r$,
\begin{equation}
\WCt_{N,n,r}(\bw_{\bM}|\bw|\bI)=\WCt_{N,s,r}(w_{m_1},\dots,w_{m_s}|\bw|\bI).
\label{eq:stabs}
\end{equation}
The normalised weight function --- and therefore every structure constant built
from it --- is \emph{independent of $n$} once $n\ge s$.  An integer containing a
free gauge parameter cannot equal a function of $(\lambda,\mu,\nu)$: on Q484's
smallest failure witness $\lambda=\mu=\emptyset$, $\nu=\{1,2\}$ at $N=2$, the
quantity $n-r$ takes \emph{every} value $\ge0$ as $n$ ranges over $n\ge2$,
including $2$ and including $0$.

\begin{thm}[the gate, Q492 closed negatively]
\label{thm:gate}
The determinant/Pfaffian stratification parameter of Motegi's boundary
specialisation is not the Q484 jump.  Precisely:
\begin{enumerate}
\item[(i)] the number of spectral parameters specialised to $1$ in
\eqref{eq:spec} is $n-s$, and $n$ is a gauge parameter removed by
\eqref{eq:stabs}, so the only gauge-invariant content of that count is
$s-r=|\bM|-|\bI|$;
\item[(ii)] $|\bM|-|\bI|$ is a cardinality gap between an \emph{evaluation
point} $\bM$ and a \emph{basis label} $\bI$ --- a (row, column) pair of the
transition matrix $\At_{\bM,\bI}$ of \S\ref{sec:setup} --- whereas
$|\nu|-\max(|\lambda|,|\mu|)$ is a gap among \emph{three basis labels}, i.e.\
three columns;
\item[(iii)] consequently the two integers are cardinality gaps in different
slots of the same matrix, and the agreement noted in the wake note (both are
``cardinality gaps between increasing sequences in $[N]$'', both vanish on the
diagonal $\bM=\bI$) is a match on one coordinate of a two-coordinate object.
\end{enumerate}
\end{thm}

\begin{rem}
This is the fourth instance this month of one error shape: a quantity that
matches on a name and on a single value, carried between two slots.  It cost
thirty minutes here because the gate was pre-registered and run at position
zero.  The distinction in (ii) is not a technicality --- it is what makes the
rest of this paper possible, because the \emph{correct} two-slot statement, the
one about rows, turns out to be the engine (\S\ref{sec:collision}).
\end{rem}

\begin{rem}[a false collision, recorded not re-litigated]
Motegi's paper says ``Laurent'' $22$ times and proves Laurent polynomiality
(L152--153, \S L3276).  That theorem is about \emph{his specialisations}
$\WCt_{N,n,r}(\bw_{\bM}|\bw|\bI)$; Q484's criterion is about the
\emph{structure constants} $c^{\nu}_{\lambda\mu}$ in that basis.  Same noun,
different object: his theorem neither implies nor contradicts mine.  The
correct relation between them, found in the course of this proof, is that his
theorem is an \emph{input} to mine --- it is exactly the regularity statement
Lemma~\ref{lem:reg} needs.  Neither support nor refutation: a hypothesis.
\end{rem}

\section{Setup: the fixed-point algebra}
\label{sec:setup}

We work in type $C$ throughout; $K^C(u)=K(u)|_{\gamma=0,\beta=1}$
(\texttt{triandc}, L374).  Let $\mathbb{F}:=\mathbb{C}(q)(w_1,\dots,w_N)$.

\begin{defn}
For $\bM,\bI\subseteq[N]$ set
$\At_{\bM,\bI}:=\WCt_{N,n,|\bI|}(\bw_{\bM}|\bw|\bI)$ for any
$n\ge\max(|\bM|,|\bI|)$.  This is well defined --- independent of $n$ --- by
\eqref{eq:stabs} for $|\bM|\ge|\bI|$ and by
\texttt{eq:typeC-stability-s-lt-r} (L2836--L2840) together with
$\At_{\bM,\bI}=0$ for $|\bM|<|\bI|$ (L2843--L2848) otherwise.  We call
$n=|\bM|$ the \emph{unpadded} and $n>|\bM|$ a \emph{padded} convention, and we
use whichever is convenient.
\end{defn}

Let $V$ be the algebra of $\mathbb{F}$-valued functions on the $2^N$
specialisation points, with pointwise product, and let
$\varphi_{\bI}\in V$ be $\bM\mapsto\At_{\bM,\bI}$.  Q484 (Lemma~2) proves
$\{\varphi_{\bI}\}$ is a basis of $V$, and we write
\begin{equation}
\varphi_{\lambda}\varphi_{\mu}=\sum_{\nu}c^{\nu}_{\lambda\mu}\varphi_{\nu},
\qquad\text{equivalently}\qquad
\At_{\bM,\lambda}\At_{\bM,\mu}=\sum_{\nu}c^{\nu}_{\lambda\mu}\At_{\bM,\nu}
\ \ \forall\bM.
\label{eq:star}
\end{equation}
Throughout, $\lambda,\mu,\nu$ denote subsets of $[N]$ and $|\cdot|$ denotes
\emph{cardinality} (Motegi's own convention: ``$r=|\bI|$ is the cardinality of
the label $\bI$'', L2790).

We use three facts about $\At$, all from Motegi.

\begin{lem}[symmetry]\label{lem:sym}
$\At_{\bM,\bI}$ depends on $\bw_{\bM}$ only through the multiset of its
entries.
\end{lem}
\begin{proof}
$F_{N,n,r}(\bu|\bw|\bI)$ is symmetric in $\bu$ (L457--L460, by the Yang--Baxter
and reflection equations); the normalisation
$\mathcal{E}^C_{N,n,r}$ (L1057--L1062) is symmetric in $\bu$ as well.
\end{proof}

\begin{lem}[triangularity]\label{lem:tri}
$\At_{\bM,\bI}=0$ unless $\bM\ge\bI$, where Motegi's order is
$\bM\ge\bI\iff|\bM|\ge|\bI|$ and $m_{s-r+j}\ge i_j$ for $1\le j\le r$.
\end{lem}
\begin{proof}
Theorem 2.6(i) of the source (\texttt{ffunctionproperties}), L625--L636; restated as
\texttt{eq:typeC-shifted-dominance}, L3250--L3258.  (Q484 measured the converse
as well, $84/84$ pairs at $N\le3$ with no violation in either direction; the
converse is \emph{not} used below except where flagged.)
\end{proof}

\begin{lem}[regularity]\label{lem:reg}
$\At_{\bM,\bI}\in\mathbb{C}(q)[w_1^{\pm1},\dots,w_N^{\pm1}]$.  Hence
$\ord_D\At_{\bM,\bI}\ge0$ for every prime divisor $D$ other than the coordinate
divisors $w_a=0$.
\end{lem}
\begin{proof}
Motegi, ``Laurent polynomiality'' for the specialisations: claimed L152--153
(``We prove support conditions, Laurent polynomiality, divisibility, and Newton
degree bounds for these specializations''), proved \S L3276ff.\ by equivariant
geometric localisation with $\lambda_y$-classes.  Deep-read first-hand.
\end{proof}

Fix a linear extension $\prec$ of Motegi's order.  By Lemmas~\ref{lem:tri}
and~\ref{lem:reg}, \eqref{eq:star} is a triangular system, solved by forward
substitution:
\begin{equation}
\At_{\nu,\nu}\,c^{\nu}_{\lambda\mu}
=\At_{\nu,\lambda}\At_{\nu,\mu}
-\sum_{\kappa\prec\nu}c^{\kappa}_{\lambda\mu}\At_{\nu,\kappa}.
\label{eq:rec}
\end{equation}
In particular $\det(\At_{\bM,\bI})=\pm\prod_{\bI}\At_{\bI,\bI}$, so
\emph{every pole of every $c^{\nu}_{\lambda\mu}$ lies on the vanishing locus of
some diagonal entry.}

\section{The four candidate divisors}
\label{sec:candidates}

Motegi evaluates the diagonal in closed form (\texttt{normalizationWC}, L1078,
written out via L1017--L1022 and L1041--L1046): with $\bI^{\mathrm c}=[N]
\setminus\bI$,
\begin{equation}
\At_{\bI,\bI}=
\underbrace{\prod_{1\le j\le k\le r}\bigl(1-w_{i_j}^{-1}w_{i_k}^{-1}\bigr)}_{\text{(L) and (S)}}
\cdot
\underbrace{\prod_{j\in\bI^{\mathrm c}}\prod_{i_l>j}\bigl(1-w_j w_{i_l}^{-1}\bigr)}_{\text{(A)}}
\cdot
\underbrace{\prod_{\substack{j\le k\\ j,k\in\bI^{\mathrm c}}}\bigl(1-qw_jw_k\bigr)
\prod_{j\in\bI^{\mathrm c}}\prod_{i_l<j}\bigl(1-qw_jw_{i_l}^{-1}\bigr)}_{\text{(Q)}} .
\label{eq:diag}
\end{equation}
This is a product over the positive roots of $C_N$, and it partitions the
candidate pole divisors into exactly four families:
\begin{center}
\begin{tabular}{@{}llll@{}}
\toprule
family & divisor & from & disposed of by \\
\midrule
(L) long root & $w_aw_b=1$, $a<b$ & $j<k$ in \eqref{eq:diag} & \S\ref{sec:long} (jump-governed) \\
(S) short root & $w_a=\pm1$ & $j=k$ in \eqref{eq:diag} & \S\ref{sec:collision} (collision) \\
(A) type $A$ root & $w_a=w_b$, $a<b$ & (A) in \eqref{eq:diag} & \S\ref{sec:collision} (collision) \\
(Q) $q$-shifted & $qw_aw_b=1$, $qw_a=w_b$ & (Q) in \eqref{eq:diag} & \S\ref{sec:qfam} (Gap~\ref{gap:q}) \\
\bottomrule
\end{tabular}
\end{center}
Naming which control covers which block is the point of the table: three of the
four families are settled unconditionally, the first is the content of the
theorem, and the last is the one honest gap.

\section{Orders of vanishing, from Motegi's branching formula}
\label{sec:orders}

Everything in \S\ref{sec:long} rests on one lemma, and Motegi's cross-sector
branching formula gives it outright.  For $\bK\subseteq\bM$ with $|\bK|=r$, put
$\bL:=\bM\setminus\bK=(l_1<\dots<l_{s-r})$.  His
\texttt{eq:typeC-one-branch-explicit} (L3203--L3236) expresses $\At_{\bM,\bI}$
as $\sum_{\bK}T_{\bK}$ with
\begin{align}
T_{\bK}=\;&(1-q)^{s-r}\!\!\prod_{\substack{j\le k\\ j,k\in\bM^{\mathrm c}}}\!\!(1-qw_jw_k)\cdot
\frac{W_{N,r}(w_{k_1},\dots,w_{k_r}|\bw|\bI)}{\prod_{a,b=1}^{r}(1-qw_{k_a}^{-1}w_{k_b})}
\notag\\
&\times
\frac{\displaystyle\prod_{1\le a\le b\le r}\!\!\bigl(1-w_{k_a}^{-1}w_{k_b}^{-1}\bigr)\;
\prod_{a=1}^{r}\prod_{b=1}^{s-r}\!\bigl(1-w_{k_a}^{-1}w_{l_b}^{-1}\bigr)}
{\displaystyle\prod_{a=1}^{r}\prod_{b=1}^{s-r}\bigl(1-w_{k_a}^{-1}w_{l_b}\bigr)}
\notag\\
&\times
\prod_{1\le a<b\le s-r}\!\!\bigl(1-qw_{l_a}^{-1}w_{l_b}^{-1}\bigr)\;
\prod_{a=1}^{s-r}\prod_{j\in\bM^{\mathrm c}}\bigl(1-qw_{l_a}^{-1}w_j\bigr).
\label{eq:branch}
\end{align}

Read off which pairs $\{x,y\}\subseteq\bM$ contribute a factor vanishing on the
\emph{long-root} divisor $w_xw_y=1$:
\begin{itemize}
\item both $x,y\in\bK$: yes, from $\prod_{a<b\le r}(1-w_{k_a}^{-1}w_{k_b}^{-1})$;
\item exactly one of $x,y$ in $\bK$: yes, from
$\prod_{a}\prod_{b}(1-w_{k_a}^{-1}w_{l_b}^{-1})$;
\item neither in $\bK$ (both in $\bL$): \textbf{no} --- such pairs appear only as
$\bigl(1-qw_{l_a}^{-1}w_{l_b}^{-1}\bigr)$, which vanishes on the $q$-shifted
divisor $qw_xw_y=1$, a different hyperplane.
\end{itemize}
No denominator in \eqref{eq:branch} vanishes on a long-root divisor: the
denominators are $\prod(1-qw_{k_a}^{-1}w_{k_b})$ and
$\prod(1-w_{k_a}^{-1}w_{l_b})$, of families (Q) and (A).  Hence:

\begin{lem}[the branch-term order]\label{lem:branchord}
Let $a<b$ with $\{a,b\}\subseteq\bM$.  Then $\ord_{\{w_aw_b=1\}}T_{\bK}\ge1$
whenever $\bK\cap\{a,b\}\ne\emptyset$.
\end{lem}

Which $\bK$ survive is also Motegi's:

\begin{lem}[admissibility]\label{lem:adm}
$T_{\bK}=0$ unless $\bK\ge\bI$, i.e.\ $k_j\ge i_j$ for all $j$.  Moreover, for
any $P\subseteq[N]$, there exists $\bK\subseteq P$ with $|\bK|=r$ and
$\bK\ge\bI$ if and only if $P\ge\bI$ in Motegi's order.
\end{lem}
\begin{proof}
The first claim is type $A$ triangularity applied to the factor
$W_{N,r}(w_{\bK}|\bw|\bI)/\prod_{a,b}(1-qw_{k_a}^{-1}w_{k_b})$: Motegi, L3263--L3269.
The second is his L3269--L3272 verbatim --- ``the necessary and sufficient
condition for $\bM$ to contain an $r$-element subset $\bK\ge\bI$ is
$\{m_{s-r+1},\dots,m_s\}\ge\bI$'' --- which is the definition of $P\ge\bI$,
since that order compares exactly the top $r$ elements of $P$ with $\bI$.
\end{proof}

\begin{lem}[Lemma A: the long-root order]\label{lem:A}
Let $a<b$ and $H_{ab}:=\{w_aw_b=1\}$.  If $\{a,b\}\subseteq\bM$ and
$\bM\setminus\{a,b\}\not\ge\bI$, then $\ord_{H_{ab}}\At_{\bM,\bI}\ge1$.
\end{lem}
\begin{proof}
By Lemma~\ref{lem:adm}, $\At_{\bM,\bI}=\sum_{\bK\ge\bI}T_{\bK}$.  Suppose some
surviving $\bK$ had $\bK\cap\{a,b\}=\emptyset$; then
$\bK\subseteq\bM\setminus\{a,b\}$ with $|\bK|=r$ and $\bK\ge\bI$, so by
Lemma~\ref{lem:adm} (second part, with $P=\bM\setminus\{a,b\}$) we would get
$\bM\setminus\{a,b\}\ge\bI$, contrary to hypothesis.  So every surviving $\bK$
meets $\{a,b\}$, and Lemma~\ref{lem:branchord} makes every surviving term
divisible by $(1-w_a^{-1}w_b^{-1})$.
\end{proof}

\begin{rem}
The same reading of \eqref{eq:branch} at $a=b$ gives the (S) analogue:
$\ord_{\{w_a=\pm1\}}\At_{\bM,\bI}\ge1$ if $a\in\bM$ and
$\bM\setminus\{a\}\not\ge\bI$, from the $a=b$ term
$(1-w_{k_a}^{-2})$.  We will not need it --- \S\ref{sec:collision} proves more
--- but it is the reason the measured order table (\S\ref{sec:verif}) has the
same shape in both families.
\end{rem}

\section{Collision divisors carry no pole, ever}
\label{sec:collision}

This section is unconditional: no hypothesis on $\lambda,\mu,\nu$.

\begin{defn}
A prime divisor $D$ is a \emph{collision divisor} with \emph{collapse map}
$\tau$ if there is a set $\mathcal{D}$ of subsets of $[N]$ and a map
$\tau:\mathcal{D}\to2^{[N]}$ such that for every $\bM\in\mathcal{D}$:
\begin{equation}
\tau(\bM)<\bM
\quad\text{and}\quad
\At_{\bM,\bI}\big|_{D}=\At_{\tau(\bM),\bI}\big|_{D}\ \ \text{for all }\bI,
\label{eq:collapse}
\end{equation}
and such that $\ord_D\At_{\bM,\bM}=0$ for $\bM\notin\mathcal{D}$.
\end{defn}

\begin{lem}[the collapses]\label{lem:collapse}
\begin{enumerate}
\item[(i)] $D=\{w_a=1\}$ is a collision divisor with
$\mathcal{D}=\{\bM:a\in\bM\}$, $\tau(\bM)=\bM\setminus\{a\}$.
\item[(ii)] $D=\{w_a=w_b\}$, $a<b$, is a collision divisor with
$\mathcal{D}=\{\bM:b\in\bM,\,a\notin\bM\}$,
$\tau(\bM)=(\bM\setminus\{b\})\cup\{a\}$.
\item[(iii)] $\At_{\bM,\bI}(-\bw)=\At_{\bM,\bI}(\bw)$; consequently
$D=\{w_a=-1\}$ is a collision divisor with the same data as in (i).
\end{enumerate}
\end{lem}
\begin{proof}
(i) Work in a padded convention with $n$ fixed, $n\ge|\bM|$.  On
$\{w_a=1\}$ the tuple $\bw_{\bM}$ has entries
$\{w_m:m\in\bM\setminus\{a\}\}$ together with $n-|\bM|+1$ copies of $1$; the
tuple $\bw_{\bM\setminus\{a\}}$, padded to the same $n$, has the same multiset.
Lemma~\ref{lem:sym} gives \eqref{eq:collapse}.  For $\tau(\bM)<\bM$: write
$a=\nu_i$ in the increasing listing of $\bM$; then
$|\bM|=|\tau\bM|+1$ and the top $|\tau\bM|$ entries of $\bM$ dominate
$\tau\bM$ entrywise, so $\bM\ge\tau\bM$, and $\ne$.

(ii) On $\{w_a=w_b\}$ with $b\in\bM$, $a\notin\bM$, the multisets of entries of
$\bw_{\bM}$ and $\bw_{(\bM\setminus\{b\})\cup\{a\}}$ agree; again
Lemma~\ref{lem:sym}.  Both sets have the same cardinality and $\tau\bM$ is
obtained from $\bM$ by decreasing one element, so $\bM>\tau\bM$.

(iii) Use the unpadded convention $n=|\bM|=s$ (legitimate, the definition is
$n$-independent; for $s<r$ both sides vanish).  Then
$(\bu,\bw)\mapsto(-\bu,-\bw)$, and we count sign changes in the bracket of
$F_{N,n,r}$ at $\beta=1$, $\gamma=0$ (L474--L487).  Invariant:
$(u_j^2-1)$, $(u_ju_k-1)$, $(u_ju_k-q)$, and each ratio
$(qu_j-u_k)/(u_j-u_k)$.  Sign-changing: $(u_j-qw_k)$ for $j\le r$,
$k=i_j+1,\dots,N$, contributing $\sum_j(N-i_j)$; $(u_j-w_k)$ for $j\le r$,
$k<i_j$, contributing $\sum_j(i_j-1)$; $(u_j-qw_k)$ for $j>r$, all $k$,
contributing $(n-r)N$; and $(\beta u_j-\gamma)=u_j$ for $j>r$, contributing
$(n-r)$.  The prefactor $\prod_k u_k^{-(N+n)}$ contributes $n(N+n)$.  Total
\[
\textstyle\sum_j(N-i_j)+\sum_j(i_j-1)+(n-r)N+(n-r)+n(N+n)
=2nN-2r+n+n^2\equiv n(n+1)\equiv0 \pmod 2 .
\]
Every factor of $\mathcal{E}^C_{N,n,r}$ --- $(1-qu_k/u_j)$, $(1-qu_jw_k)$,
$(1-qu_ju_k)$, $(1-qw_jw_k)$ --- is invariant.  So $\At$ is invariant.  Now
$\At_{\bM,\bI}|_{w_a=-1}$ equals $\At_{\bM,\bI}(\bv)|_{v_a=1}$ under
$\bw=-\bv$, which by (i) equals $\At_{\bM\setminus\{a\},\bI}(\bv)|_{v_a=1}
=\At_{\bM\setminus\{a\},\bI}|_{w_a=-1}$.
\end{proof}

\begin{thm}[collision divisors]\label{thm:collision}
Let $D$ be a collision divisor.  Then $\ord_D c^{\nu}_{\lambda\mu}\ge0$ for all
$\lambda,\mu,\nu$.  In particular no $c^{\nu}_{\lambda\mu}$ has a pole on
$w_a=1$, on $w_a=-1$, or on $w_a=w_b$.
\end{thm}
\begin{proof}
Induction along $\prec$.  Assume $\ord_Dc^{\kappa}_{\lambda\mu}\ge0$ for all
$\kappa\prec\nu$ (vacuous at the start).  By Lemma~\ref{lem:reg} every $\At$ is
regular on $D$.

If $\nu\notin\mathcal{D}$ then $\ord_D\At_{\nu,\nu}=0$ and \eqref{eq:rec} gives
$\ord_Dc^{\nu}\ge0$ immediately.

So let $\nu\in\mathcal{D}$ and put $\nu':=\tau(\nu)<\nu$.  Since
$|\nu'|<|\nu|$ or $\nu'\ne\nu$ with $\nu>\nu'$, triangularity gives
$\At_{\nu',\nu}=0$, hence by \eqref{eq:collapse}
$\ord_D\At_{\nu,\nu}\ge1$.  We show the right side of \eqref{eq:rec} vanishes on
$D$ to at least the same order; it suffices to show it vanishes on $D$, since
$\ord_D\At_{\nu,\nu}$ will be matched after dividing --- precisely, we show
\begin{equation}
\Bigl[\At_{\nu,\lambda}\At_{\nu,\mu}-\sum_{\kappa\prec\nu}c^{\kappa}\At_{\nu,\kappa}\Bigr]
\Big|_D=0 .
\label{eq:numvanish}
\end{equation}
Apply \eqref{eq:collapse} to each $\At_{\nu,\bullet}$:
the left side of \eqref{eq:numvanish} equals
$\bigl[\At_{\nu',\lambda}\At_{\nu',\mu}-\sum_{\kappa\prec\nu}c^{\kappa}
\At_{\nu',\kappa}\bigr]|_D$, all terms being regular on $D$ by
Lemma~\ref{lem:reg} and the inductive hypothesis.  Now evaluate \eqref{eq:star}
at $\bM=\nu'$:
\[
\At_{\nu',\lambda}\At_{\nu',\mu}=\sum_{\kappa}c^{\kappa}\At_{\nu',\kappa},
\]
where by Lemma~\ref{lem:tri} the sum is over $\kappa$ with $\nu'\ge\kappa$.
Every such $\kappa$ satisfies $\kappa\le\nu'<\nu$, hence $\kappa\prec\nu$; and
every $\kappa\prec\nu$ with $\nu'\not\ge\kappa$ has $\At_{\nu',\kappa}=0$.  So
the two sums agree term by term, and \eqref{eq:numvanish} holds.

Finally, \eqref{eq:numvanish} says the numerator lies in the prime ideal of
$D$, so $\ord_D(\text{numerator})\ge1$.  It remains to check
$\ord_D\At_{\nu,\nu}=1$, i.e.\ that the collision divisor occurs in the
diagonal entry to multiplicity exactly one.  This is visible in
\eqref{eq:diag}: for $D=\{w_a=\varepsilon\}$ the only factor of
$\At_{\nu,\nu}$ vanishing on $D$ is the single $j=k=a$ term
$(1-w_a^{-2})=-w_a^{-2}(w_a-1)(w_a+1)$, in which $(w_a-\varepsilon)$ occurs once;
for $D=\{w_a=w_b\}$ it is the single factor $(1-w_aw_b^{-1})$ with
$a\in\nu^{\mathrm c}$, $b\in\nu$, $b>a$.  All factors listed in
\eqref{eq:diag} are pairwise non-associate irreducibles, so no multiplicity
accumulates.  Hence $\ord_Dc^{\nu}\ge1-1=0$.
\end{proof}

\begin{rem}[where this comes from, physically]
$[K^C(u)]^1_0=1-u^{-2}$ vanishes at $u=\pm1$, which is why a row carrying
$u=\pm1$ is forced and removable --- that is Motegi's own engine (L1750,
L2094), and \eqref{eq:stab} is its normalised form.  The brief asked me to name
the component I skip: $1-u^{-2}$ vanishes at \emph{both} $u=1$ and $u=-1$, and
Motegi only ever uses $u=+1$.  Lemma~\ref{lem:collapse}(iii) answers it.  The
$u=-1$ branch is \textbf{not} trivial --- it is a genuine second collision
divisor, with the same collapse map --- and it is reached by the symmetry
$\bw\mapsto-\bw$, not by Motegi's argument.  Unlike the $q=-1$ collapse of
2026-10-15, there is mathematics on this component.
\end{rem}

\section{The long-root divisors: the jump appears}
\label{sec:long}

The long-root divisors $w_aw_b=1$ are \emph{not} collision divisors: setting
$w_b=w_a^{-1}$ does not make $\bw_{\bM}$ coincide with any other specialisation
point, because $w_a^{-1}$ is not among the $w_c$ or the padding $1$'s.  This is
the whole asymmetry, and it is where poles live.

\begin{lem}[the combinatorial step]\label{lem:comb}
Let $\nu\ge\lambda$ and $\nu\ge\mu$.  The following are equivalent:
\begin{enumerate}
\item[(a)] there is a pair $a<b$ with $\{a,b\}\subseteq\nu$,
$\nu\setminus\{a,b\}\ge\lambda$ and $\nu\setminus\{a,b\}\ge\mu$;
\item[(b)] $|\nu|-\max(|\lambda|,|\mu|)\ge2$.
\end{enumerate}
\end{lem}
\begin{proof}
(a)$\Rightarrow$(b): $\nu\setminus\{a,b\}\ge\lambda$ forces
$|\nu|-2\ge|\lambda|$, and likewise $|\nu|-2\ge|\mu|$.

(b)$\Rightarrow$(a): let $\{a,b\}$ be the two \emph{smallest} elements of
$\nu$ (available since $|\nu|\ge2$), and $\nu':=\nu\setminus\{a,b\}$, so
$|\nu'|=|\nu|-2\ge\max(|\lambda|,|\mu|)$.  Motegi's order compares only the
top $|\lambda|$ elements of the larger set with $\lambda$; since
$|\nu'|\ge|\lambda|$, the top $|\lambda|$ elements of $\nu'$ are the top
$|\lambda|$ elements of $\nu$.  Hence $\nu\ge\lambda$ gives $\nu'\ge\lambda$,
and similarly $\nu'\ge\mu$.
\end{proof}

\begin{rem}
This is where the $\max$ in Q484's criterion comes from, and why it looked
strange.  The criterion reads ``$\min(|\nu|-|\lambda|,|\nu|-|\mu|)\le1$'' ---
an \emph{or} --- and the reason is that a pole needs the deleted pair to stay
above \emph{both} $\lambda$ and $\mu$, an \emph{and} whose negation is an
\emph{or}.  The two deleted elements are the two smallest of $\nu$; the cost
``$2$'' is the size of a long root's support.
\end{rem}

\begin{thm}[main: jump $\le1$ implies no long-root pole]\label{thm:main}
Suppose $|\nu|-\max(|\lambda|,|\mu|)\le1$.  Then
$\ord_{H_{ab}}c^{\nu}_{\lambda\mu}\ge0$ for every $a<b$.  Combined with
Theorem~\ref{thm:collision} and Gap~\ref{gap:q}, $c^{\nu}_{\lambda\mu}$ is then
a Laurent polynomial in $w_1,\dots,w_N$ over $\mathbb{C}(q)$.
\end{thm}
\begin{proof}
Write $m:=\max(|\lambda|,|\mu|)$ and assume without loss of generality
$|\lambda|=m$.  Fix $a<b$ and put $H:=H_{ab}$.  We induct along $\prec$, the
claim at $\nu$ being $\ord_Hc^{\nu}\ge0$ for all $\nu$ with $|\nu|\le m+1$.

If $\{a,b\}\not\subseteq\nu$ then $\ord_H\At_{\nu,\nu}=0$ by \eqref{eq:diag}
--- the factor $(1-w_a^{-1}w_b^{-1})$ occurs there only for $a,b\in\nu$ --- and
\eqref{eq:rec} with Lemma~\ref{lem:reg} and the inductive hypothesis gives
$\ord_Hc^{\nu}\ge0$.

So let $\{a,b\}\subseteq\nu$ and $\nu':=\nu\setminus\{a,b\}$.  Since
$|\nu'|=|\nu|-2<|\nu|$, triangularity gives $\nu'\not\ge\nu$, so
Lemma~\ref{lem:A} applies with $\bM=\bI=\nu$ and
$\ord_H\At_{\nu,\nu}\ge1$; by \eqref{eq:diag} the factor occurs to order
exactly $1$, so $\ord_H\At_{\nu,\nu}=1$.  It therefore suffices to show the
right side of \eqref{eq:rec} has $\ord_H\ge1$, and we do it termwise.

\emph{Term $\At_{\nu,\lambda}\At_{\nu,\mu}$.}  We have
$|\nu'|=|\nu|-2\le m-1<m=|\lambda|$, so $\nu'\not\ge\lambda$ for cardinality
reasons.  Lemma~\ref{lem:A} gives $\ord_H\At_{\nu,\lambda}\ge1$, and
$\ord_H\At_{\nu,\mu}\ge0$ by Lemma~\ref{lem:reg}.

\emph{Terms $c^{\kappa}\At_{\nu,\kappa}$, $\kappa\prec\nu$.}  Only $\kappa$
with $\At_{\nu,\kappa}\ne0$ matter, so $\nu\ge\kappa$, hence
$|\kappa|\le|\nu|\le m+1$ and the inductive hypothesis gives
$\ord_Hc^{\kappa}\ge0$.  Only $\kappa$ with $c^{\kappa}\ne0$ matter, so by the
support bound (Q484, Theorem~A) $\kappa\ge\lambda$ and $\kappa\ge\mu$, hence
$|\kappa|\ge m$.  Then $|\nu'|=|\nu|-2\le m-1<m\le|\kappa|$, so
$\nu'\not\ge\kappa$ and Lemma~\ref{lem:A} gives
$\ord_H\At_{\nu,\kappa}\ge1$.  Therefore
$\ord_H(c^{\kappa}\At_{\nu,\kappa})\ge0+1=1$.

Both terms have $\ord_H\ge1=\ord_H\At_{\nu,\nu}$, so
$\ord_Hc^{\nu}\ge0$.
\end{proof}

\begin{rem}[why the induction closes]
The step that makes this work is that the hypothesis is inherited: every
$\kappa$ entering the recursion has $|\kappa|\le|\nu|\le
\max(|\lambda|,|\mu|)+1$, so $\kappa$ also has jump $\le1$ and the inductive
hypothesis is available.  Equivalently: \emph{jump is monotone in $|\nu|$ along
Motegi's order}, which is why a pole cannot be created below $\nu$ and
inherited by a $\nu$ of small jump.
\end{rem}

\subsection{The exact pole divisor: creation and propagation}

The mechanism says more than Theorem~\ref{thm:main}.  Reading
Lemma~\ref{lem:comb} locally at each $\kappa$ rather than at $\nu$ gives a
prediction of the \emph{divisor}, not just of Laurent-ness.

\begin{conj}[exact pole divisor; verified, see \S\ref{sec:verif}]
\label{conj:divisor}
For all $\lambda,\mu,\nu$ the set of long-root divisors on which
$c^{\nu}_{\lambda\mu}$ has a pole is
\[
\bigl\{H_{ab}\;:\;\exists\,\kappa\ \text{with}\ \nu\ge\kappa\ge\lambda,\
\kappa\ge\mu,\ \{a,b\}\subseteq\kappa,\
\kappa\setminus\{a,b\}\ge\lambda,\ \kappa\setminus\{a,b\}\ge\mu\bigr\}.
\]
\end{conj}

A pole is \emph{created} at $\kappa$ when the bracket condition holds there, and
then \emph{propagates} upward: for $\nu>\kappa$ the term
$c^{\kappa}\At_{\nu,\kappa}$ in \eqref{eq:rec} carries it, and
$\At_{\nu,\kappa}$ need not vanish on $H_{ab}$.  This is why, e.g.,
$c^{\{1,3\}}_{\emptyset\emptyset}$ at $N=3$ has a pole on $w_1w_2=1$ although
$2\notin\{1,3\}$: the pole is created at $\kappa=\{1,2\}\le\{1,3\}$.  I had
first predicted only the created poles and the measurement corrected me;
Conjecture~\ref{conj:divisor} is the corrected statement.  Note that it is
consistent with Theorem~\ref{thm:main}: if $\nu$ has jump $\le1$ then every
$\kappa\le\nu$ in the support has $|\kappa|\le|\nu|$, hence jump $\le1$, hence
by Lemma~\ref{lem:comb} no $\kappa$ satisfies the bracket.

\section{The $q$-shifted family}
\label{sec:qfam}

\begin{gapthm}[the $q$-shifted divisors]\label{gap:q}
For $D$ in family (Q) --- $qw_aw_b=1$ or $qw_a=w_b$ --- I verify
computationally (\S\ref{sec:verif}) that
\begin{equation}
\ord_D\At_{\nu,\nu}\le\ord_D\At_{\nu,\bI}\qquad\text{for all }\bI\le\nu,
\label{eq:qrow}
\end{equation}
which closes the case at once: with $e:=\ord_D\At_{\nu,\nu}$, \eqref{eq:rec}
gives $\ord_D(\text{numerator})\ge\min(2e,0+e)=e$ and so
$\ord_Dc^{\nu}\ge0$ by the same induction as before.  I do \emph{not} prove
\eqref{eq:qrow}.  What I can prove is the structural half: the contribution of
the normalisation $\mathcal{E}^C$ to $\ord_D\At_{\bM,\bI}$ depends on $\bM$
alone and not on $\bI$ --- its denominator
$\prod_{j\le k}(1-qw_jw_k)\prod_{j<k\in\bM}(1-qw_{m_j}w_{m_k})$ and numerator
$\prod(1-qw_{m_k}/w_{m_j})\prod(1-qw_{m_j}w_k)$ are functions of $\bM$ only ---
so that \eqref{eq:qrow} is a statement about the $\bI$-dependence of the
$W_{N,r}$ and bottom-boundary blocks of \eqref{eq:branch} alone.  Status:
verified at $N=2,3$ over all $\bI\le\nu$, $9$ divisors at $N=2$ and $9$ at
$N=3$, no violation; not proved.
\end{gapthm}

\begin{gapthm}[non-vanishing, the converse direction]\label{gap:nonvanishing}
Theorem~\ref{thm:main} is the ``jump $\le1\Rightarrow$ Laurent'' half.  The
converse --- jump $\ge2$ implies a pole \emph{actually occurs} --- requires that
the created pole have nonzero residue, i.e.\ that the numerator of
\eqref{eq:rec} does \emph{not} vanish on $H_{ab}$ when Lemma~\ref{lem:comb}(a)
holds.  Lemma~\ref{lem:A} is a divisibility statement and gives only lower
bounds on orders, so it cannot supply this.  Equivalently, one needs the
converse of Lemma~\ref{lem:tri} ($\bM\ge\bI\Rightarrow\At_{\bM,\bI}\ne0$,
measured $84/84$ in Q484) together with the absence of conspiratorial
cancellation among the surviving $\bK$-terms of \eqref{eq:branch}.
Status: measured.  $225/225$ triples at $N=2,3$ in Q484 with zero error in
either direction; and, much more sharply, Conjecture~\ref{conj:divisor} holds
exactly, $20/20$ and $115/115$ (\S\ref{sec:verif}).
\end{gapthm}

\section{Computational verification, with controls and plants}
\label{sec:verif}

All computation is pure Python $+$ SymPy over
$\mathbb{C}(q)(w_1,\dots,w_N)$ (SageMath is not installed in this container).
The engine is the Q484 engine \texttt{q484\_engine.py}, which evaluates
$F_{N,n,r}(\bu|\bw|\bI)$ by brute-force summation over $S_n$ directly from the
symmetrisation formula (L474--L487) --- i.e.\ it does \emph{not} use the
branching formula \eqref{eq:branch} on which \S\ref{sec:orders} rests, so the
two are independent.  Scripts: \texttt{proofs/code-1016-q492/}.

\paragraph{Reproduction of Q484.}
At $N=2$, $29/30$ triples Laurent, the unique failure
$c^{\{1,2\}}_{\emptyset\emptyset}$ with pole divisor $(w_1w_2-1)$.  At $N=3$,
$102/115$ Laurent on the $\mu\ge\lambda$ half, i.e.\ $176/195$ in full.  Jump
distribution: jump $0$ gives $58+12$ Laurent and $0$ poles; jump $1$ gives
$44+7$ Laurent and $0$ poles; jump $2$ gives $0$ Laurent and $12+1$ poles; jump
$3$ gives $0$ Laurent and $1$ pole.

\paragraph{Control 1: the exact pole divisor (independent mechanism).}
Conjecture~\ref{conj:divisor} is pure combinatorics on subsets; the measured
divisor comes from SymPy rational arithmetic and \texttt{factor\_list}.  The two
agree on \textbf{$20/20$} nonzero constants at $N=2$ and \textbf{$115/115$} at
$N=3$.  This is a far stronger test than the binary criterion: it checks
\emph{which} hyperplanes, with no free parameter.  For instance at $N=3$,
$c^{\{1,2,3\}}_{\{1\}\{2\}}$ is predicted to have poles on $w_1w_2=1$ and
$w_1w_3=1$ but \emph{not} on $w_2w_3=1$, and the measured denominator is
$(w_1w_2-1)(w_1w_3-1)$.

\paragraph{Control 2: no pole on a collision or $q$-divisor.}
Theorem~\ref{thm:collision} predicts that no denominator ever contains a factor
from families (S), (A) or (Q).  Measured: \textbf{$20/20$} at $N=2$ and
\textbf{$115/115$} at $N=3$ --- every pole divisor observed, across all
triples, is a product of long roots $w_aw_b-1$ only.

\paragraph{Order tables (Lemma~\ref{lem:A} and its (S), (A) analogues).}
Measured as exact equalities, not just the lower bounds the proof uses:
$\ord_{H_{ab}}\At_{\bM,\bI}=[\{a,b\}\subseteq\bM]\cdot
[\bM\setminus\{a,b\}\not\ge\bI]$, confirmed on $10$ pairs $(\bM,\bI)$ at $N=2$
and $105$ at $N=3$;
$\ord_{\{w_a=\varepsilon\}}\At_{\bM,\bI}=[a\in\bM][\bM\setminus\{a\}\not\ge\bI]$
for both $\varepsilon=\pm1$, confirmed at $N=2,3$;
$\ord_{\{w_a=w_b\}}\At_{\bM,\bI}=[a\notin\bM][b\in\bM]
[(\bM\setminus\{b\})\cup\{a\}\not\ge\bI]$, confirmed at $N=2,3$ --- note this
last is exactly the collapse map of Lemma~\ref{lem:collapse}(ii), measured
independently of it.
Also verified, since Theorem~\ref{thm:collision} needs it: on every collision
divisor $D$ and for every $\bI$, $\ord_D\At_{\bI,\bI}\in\{0,1\}$ --- the
multiplicity-one claim at the end of that proof --- at $N=2$ ($6$ divisors) and
$N=3$ ($9$ divisors), no value $\ge2$ occurring.

\paragraph{Control 3: the $\bw\mapsto-\bw$ symmetry.}
Lemma~\ref{lem:collapse}(iii) asserts $\At(-\bw)=\At(\bw)$.  Verified on
\textbf{$10/10$} nonzero entries at $N=2$ and \textbf{$35/35$} at $N=3$.  This
is the control for the $u=-1$ component, and it is \emph{not} over an inherited
domain: the entries are the same ones the order tables run on.

\paragraph{Plant P1 (K-matrix off-diagonal): FIRED.}
Replace the vanishing factor $(u_j^2-1)\mapsto(u_j^2-q)$ in the $j\le r$ block,
which moves the short-root collision locus from $w_a=\pm1$ to $w_a^2=q$ --- a
locus that is \emph{not} a collision of specialisation points.
Theorem~\ref{thm:collision} should then fail.  It does: Control 2 drops to
$12/20$ and the new denominators are exactly $q-w_a^2$.  The plant fires, and it
fires \emph{on the predicted family}, which is the informative part.

\paragraph{Plant P2 (R-matrix weight): FIRED.}
Replace $(u_ju_k-1)\mapsto(u_ju_k-q)$ in the $1\le j<k\le r$ block, moving the
long-root divisor to $qw_aw_b=1$.  Control 1 drops to $19/20$ --- the
$\{1,2\}$ pole leaves $w_1w_2=1$, exactly as predicted --- and Control 2 to
$14/20$, with new denominators $q-w_aw_b$.  Both plants fire.

\paragraph{The stratum each control validates, and what it cannot see.}
Following the 2026-10-15 fault in which a $14/14$ diagonal control was blind to
a plant because $\bM=\bI$ empties the index range $r+1\le j<k\le n$: the
controls here run on the \emph{off-diagonal} stratum by construction, since
Control 1 and Control 2 are statements about $c^{\nu}_{\lambda\mu}$ for all
triples, including $|\nu|>|\lambda|,|\mu|$, where $s>r$ and all blocks of
\eqref{eq:branch} are live.  The order tables run over all $(\bM,\bI)$ with
$\At\ne0$, which at $N=3$ is $35$ pairs of which $27$ are off-diagonal.  What
the controls do \emph{not} cover: $N\ge4$ (Lemma~\ref{lem:comb}'s content is
$N$-independent, but Conjecture~\ref{conj:divisor} is verified only at
$N\le3$), and type $B$ (every number here is type $C$; see
\S\ref{sec:open}).

\section{Open threads}
\label{sec:open}

\begin{enumerate}
\item \textbf{Gap~\ref{gap:nonvanishing} (the converse).}  Proving
jump $\ge2\Rightarrow$ pole needs the converse of Lemma~\ref{lem:tri} plus
non-cancellation among the $\bK$-terms of \eqref{eq:branch}.  The natural route
is to show the leading coefficient on $H_{ab}$ is a sum of \emph{like-signed}
contributions; Motegi's branching formula \eqref{eq:branch} makes each
$T_{\bK}$ explicit, so this is a finite computation per $\bK$-class and not
obviously hard.  This is the next prove session.

\item \textbf{Gap~\ref{gap:q} (the $q$-family).}  Prove \eqref{eq:qrow}.  Half
the statement is reduced to the $\bI$-dependence of the $W_{N,r}$ block.

\item \textbf{Type $B$.}  Every number in this paper and in Q484 is type $C$.
Nothing in \S\S\ref{sec:collision}--\ref{sec:long} uses $\gamma=0$ except
Lemma~\ref{lem:collapse}(iii), whose sign count would have to be redone for
$WB$ (there $\gamma=-1$, and the extra prefactor
$\prod_k(1+u_k^{-1})^{-1}$ is \emph{not} invariant under
$\bu\mapsto-\bu$ --- it has a pole at $u_k=-1$).  So the $u=-1$ collision
divisor is exactly where type $B$ may differ from type $C$, which is a sharper
prediction than Q484's open thread 3 and is cheap to test.  Motegi defines
type $B$ $\bM$-smallness at L1346 in the same file.

\item \textbf{$N\ge4$.}  Conjecture~\ref{conj:divisor} is verified only at
$N\le3$.  Lemma~\ref{lem:comb} and Theorem~\ref{thm:main} are $N$-independent.

\item \textbf{What Q492 was really asking.}  Theorem~\ref{thm:gate} kills the
identification of the jump with Motegi's $n-s$.  But \S\ref{sec:collision}
supplies the structural statement the question wanted, in the slot the gate
says is the right one: a puzzle-style rule fails not because the object is ``a
Pfaffian read in determinant coordinates'' but because \emph{the long-root
divisors are the ones with no collision to collapse them}.  Whether that is
related to the determinant/Pfaffian dichotomy of \texttt{1806.00933} remains
unknown to me: that paper is at extraction level \emph{abstract} and is not on
disk.  Nothing above depends on it.
\end{enumerate}

\section*{Source hygiene}

Of the sources this paper composes, exactly one is read first-hand, and it is
the one every proof step cites: Motegi, arXiv:2610.15344, e-print source
\texttt{typebcweightfunctions\_arxivversionfinal-3.tex}, $216{,}486$ B, $6430$
lines, deep-read.  Every line number above was read in that file during this
session.  The Pfaffian side of the original dichotomy
(\texttt{1806.00933}, \texttt{2309.08446}, \texttt{2603.20865},
\texttt{math/0008184}) is \textbf{unread} --- abstract or title-only, none on
disk --- and consequently \emph{no statement in this paper rests on it}; the
word Pfaffian appears here only in \S\ref{sec:gate} and
\S\ref{sec:open} item 5, in both cases to say that the connection is not
established.  Q484's results (support bound, triangularity converse,
the $225$-triple measurement) are my own, graded in
\texttt{proofs/registry/}, and are cited as such.

\end{document}
